\section{Related work}
\label{sec:related}
\paragraph{MRI biomarkers and learned pCR prediction.}
MRI-based response assessment has a substantial clinical and quantitative foundation. Functional tumor volume and related measurements have been studied in longitudinal breast MRI cohorts, including I-SPY \citep{hylton2012mri,hylton2016ftv}. The ACRIN 6698 study additionally examined diffusion-derived measurements in relation to response \citep{ACRIN6698diffusion2018}. These studies motivate imaging-based monitoring but do not establish that a learned model using different inputs or a different evaluation split has equivalent clinical validity. Our experiments use selected DCE phases and baseline clinical variables; diffusion imaging is not an input to the present model.

Learning-based approaches can combine complementary image sequences or repeated examinations. Synthesis-based representation learning has been explored for multi-sequence 3D/4D MRI \citep{han2024synthesis}, and recent studies have examined spatiotemporal or early-treatment information for pCR prediction \citep{tang2025spatiotemporal,du2025early}. Such studies provide relevant methodological context, but differences in patient selection, imaging availability, endpoints, and holdout construction prevent direct numerical comparisons with our results. We therefore do not claim superiority to published systems from cross-paper metric comparisons. Reimplementation of suitable external baselines under the same patient split remains necessary.

\paragraph{Predictive latent representations.}
Predicting a target embedding rather than reconstructing the input image is a central idea in joint-embedding predictive learning \citep{jepa2023}. We use an exponential-moving-average target encoder to define latent states and supervise autonomous future predictions. The present objective is a supervised pCR loss combined with a masked multi-step latent-state loss. It does not activate a separate single-step teacher-forcing JEPA loss, image reconstruction loss, or latent-variance regularizer. This objective-level distinction matters when interpreting the internal experiment name \nojepa{}: that control removes future-state supervision, while retaining the common fusion architecture.

\paragraph{Attention and volumetric image representation.}
Transformer attention provides a flexible way to combine image tokens and contextual information \citep{vaswani2017}. Compact convolutional representations can reduce the cost of processing volumetric MRI, and vector-quantized generative models provide one source of pretrained convolutional encoders \citep{vqgan2021}. Here the pretrained codec contributes a continuous convolutional encoder. The quantizer and image decoder are bypassed, and the main experiments do not synthesize MRI volumes. The final attention weights summarize latent tokens that have already undergone spatial mixing; they are therefore not validated lesion segmentations or anatomical explanations.
